% !TeX root = ../main.tex
\section{Numerical examples and independent checks}
\label{sec:experiments}

We use synthetic moduli to illustrate the directional pressure response.
The examples also check the finite-deformation calculation. Stresses and stiffnesses use the same
reference stress unit \(K_*\). Let \(\phi_{s0}=0.6\),
\(\mathcal{K}=7K_*\), and prescribe the mineral stiffness in an orthonormal Mandel
basis ordered as \(11,22,33,23,13,12\):
\begin{equation}
 \frac{\mathbb{C}_s}{K_*}=\begin{pmatrix}
 50&12&10&0&0&0\\
 12&60&14&0&0&0\\
 10&14&70&0&0&0\\
 0&0&0&20&0&0\\
 0&0&0&0&24&0\\
 0&0&0&0&0&28
 \end{pmatrix}.
 \label{eq:example-mineral-stiffness}
\end{equation}
The shear strain coordinates in this basis are \(\sqrt2\varepsilon_{ij}\).
The mineral is positive definite and \(\mathcal{K}_s=28K_*\), so
\(0<\mathcal{K}<\phi_{s0}\mathcal{K}_s\). Equation \eqref{eq:drained-stiffness-restriction}
determines the complete drained stiffness, including its shear entries.
The reference Biot components are \(0.7000\), \(0.7583\), and \(0.7917\).

The isotropic comparison uses mineral bulk modulus \(K_s=28K_*\) and
drained bulk modulus \(K=7K_*\), equal to the corresponding spherical-strain
moduli of the anisotropic material, and the same reference solid fraction. Its
mineral shear modulus is \(16.8K_*\), the mean of the five deviatoric
stiffness modes of \(\mathbb C_s\), divided by two. The comparison therefore
preserves the mineral spherical response and its average deviatoric
stiffness, while removing directional dependence. Both materials obey the
drained compliance restriction; their drained stiffnesses are calculated,
not selected independently. These are homogeneous constitutive
calculations with synthetic parameters, rather than finite-element
simulations or calibration to a measured material.

\subsection{Pressure change at fixed deformation}
Hold \(\mathbf F=\mathbf I\) and increase \(p/K_*\) from zero to six.
The mineral equation is identical for the two materials on this path,
because the skeleton shape term vanishes and the mineral moduli
\(\mathcal{K}_s\) and \(K_s\) have the same value.
Consequently their mineral volume and solid fraction coincide. Their
pressure-induced stresses differ: the anisotropic mineral gives three
normal coupling coefficients, while the isotropic mineral gives a single
coefficient. Figure~\ref{fig:spherical-pressure} shows the coefficients,
total stresses, and phase volumes. The final panel measures the error
from replacing the pressure integral in
\eqref{eq:integrated-pressure-response} by the instantaneous product
\(-p\mathbf B(\mathbf F,p)\). This difference grows as pressure changes
mineral volume and hence the coupling tensor.
\begin{figure}[htbp]
 \centering
 \resizebox{.91\textwidth}{!}{\import{build/conformal/}{pressure_response.pgf}}
 \caption{Fixed-deformation pressure experiment, \(\mathbf F=\mathbf I\).
 All stresses and pressures are expressed in \(K_*\). The curves use 121
 pressure states on \(0\le p/K_*\le6\). The anisotropic and isotropic
 minerals have the same mineral-volume curve but different directional
 stress sensitivities. Panel (d) compares the actual total stress with
 \(-p\mathbf B(p)\); the drained stress is zero on this path.}
 \label{fig:spherical-pressure}
\end{figure}

\subsection{Shape change and directional pressure sensitivity}
To separate shape effects from skeleton volume, prescribe
\begin{equation}
 \mathbf F(\gamma)=
 \begin{pmatrix}
 e^{0.12}&\gamma&0\\
 0&e^{-0.04}&0\\
 0&0&e^{-0.08}
 \end{pmatrix},\qquad J=1,\qquad p=2K_*.
 \label{eq:experiment-shear-gradient}
\end{equation}
The diagonal stretches are fixed while \(-0.8\le\gamma\le0.8\).
Figure~\ref{fig:spherical-shear} shows that the anisotropic mineral changes
volume as skeleton shape changes, even though \(J\) and pressure remain
fixed. Its off-diagonal coefficient \(B_{12}\) is nonzero away from
\(\gamma=0\). A pressure increment at fixed deformation therefore changes
shear stress as well as normal stress. For the isotropic mineral, the
mineral-volume equation depends only on \(J\) and \(p\); its volume and
spherical Biot tensor are constant along this isochoric path.
\begin{figure}[htbp]
 \centering
 \resizebox{\textwidth}{!}{\import{build/conformal/}{shear_response.pgf}}
 \caption{Isochoric shape-change experiment at \(p=2K_*\), using
 \eqref{eq:experiment-shear-gradient} and 161 equally spaced values of
 \(\gamma\). Mineral volume--shape coupling changes \(\bar J\) and
 generates a shear pressure sensitivity in the anisotropic material.
 The isotropic comparison has \(B_{12}=0\).}
 \label{fig:spherical-shear}
\end{figure}

The Biot tensor also determines how a pressure change alters the traction
on a plane. For current unit normal \(\mathbf n\) and unit tangent
\(\mathbf e_t\), define
\begin{equation}
 \mathbf n=(\cos\theta,\sin\theta,0)^T,\qquad
 \mathbf e_t=(-\sin\theta,\cos\theta,0)^T.
 \label{eq:experiment-plane-directions}
\end{equation}
Equation \eqref{eq:cauchy-pressure-tangent} gives the negative pressure
derivatives of normal and tangential traction as
\(\mathbf n\cdot\mathbf B\mathbf n\) and
\(\mathbf e_t\cdot\mathbf B\mathbf n\), respectively.
Figure~\ref{fig:spherical-directional} compares the undeformed skeleton
with \eqref{eq:experiment-shear-gradient} at \(\gamma=0.65\), both at
\(p=2K_*\). Anisotropy is visible both in the normal sensitivity and in
the nonzero tangential sensitivity. Changing the deformation changes the
preferred directions. The isotropic comparison gives a circular normal
response and zero tangential response at every angle.
\begin{figure}[htbp]
 \centering
 \resizebox{.93\textwidth}{!}{\import{build/conformal/}{directional_response.pgf}}
 \caption{Directional pressure sensitivities at \(p=2K_*\), sampled at
 one-degree intervals. The finite-shear state is
 \eqref{eq:experiment-shear-gradient} with \(\gamma=0.65\). Directions
 refer to fixed current spatial axes, and both plotted quantities are
 dimensionless. Tangential pressure sensitivity is a direct consequence
 of a nonspherical Biot tensor.}
 \label{fig:spherical-directional}
\end{figure}

\subsection{Internal and physical rotations}
We consider two different rotation experiments. First, hold the physical
state \(\mathbf F=\mathbf F(0.65)\), \(p=2K_*\) fixed and vary
\(\mathbf R_A\) through rotations about the third axis. Second, hold
\(\mathbf R_A=\mathbf I\) and rotate the entire deformation by the
same angle, \(\mathbf F\mapsto\mathbf Q\mathbf F\).
Figure~\ref{fig:spherical-rotation} shows constant components in the first
experiment and the usual spatial tensor transformation in the second.
The true-frame mineral stress changes in the first experiment, but its
mixture-frame representation, energy, mineral volume, and Biot tensor
remain unchanged. The distinction is between changing the internal
representation and rotating the physical body; neither operation changes
the prescribed material stiffness in its reference axes.
\begin{figure}[htbp]
 \centering
 \resizebox{.93\textwidth}{!}{\import{build/conformal/}{rotation_response.pgf}}
 \caption{Internal-frame change versus superposed physical rotation.
 The base state is \(\mathbf F(0.65)\), \(p=2K_*\); 121 angles span
 zero to \(180^\circ\). At fixed \(\mathbf F\), varying
 \(\mathbf R_A\) leaves the response unchanged. Rotating the body gives
 \(\mathbf B\mapsto\mathbf Q\mathbf B\mathbf Q^T\).}
 \label{fig:spherical-rotation}
\end{figure}

\subsection{Laterally constrained extension}
For a homogeneous layer whose lateral stretches remain one and whose
upper surface has zero total normal stress, set
\(\mathbf F=\operatorname{diag}(1,1,\exp\varepsilon_{33})\) and solve
\(\sigma_{33}=0\) at each pressure. This scalar boundary-condition problem
is coupled to the mineral-volume equation. Figure~\ref{fig:spherical-layer}
shows axial extension, lateral reaction stresses, and mineral volume.
The anisotropic material has distinct lateral reactions and an axial
response different from the isotropic comparison. These are boundary
condition effects as well as constitutive effects: lateral deformation is
prevented while axial deformation is free. The calculation illustrates how
a tensorial pressure response enters a measurable constrained deformation.
\begin{figure}[htbp]
 \centering
 \resizebox{\textwidth}{!}{\import{build/conformal/}{constrained_layer.pgf}}
 \caption{Homogeneous laterally constrained layer with
 \(F_{11}=F_{22}=1\) and \(\sigma_{33}=0\). Pressure is sampled at 121
 points on \(0\le p/K_*\le6\); stresses are expressed in \(K_*\).
 The axial stretch is \(\lambda_3=\exp\varepsilon_{33}\). Differences
 between the lateral reactions reflect the anisotropic mineral law.}
 \label{fig:spherical-layer}
\end{figure}
